\SourceLeaf{154}
One of the principal reasons I later set aside the metaphysics used in the Theoria Motus Corporum Coelestium for the method of least squares is one I have not mentioned publicly. I must in every respect consider it less important to determine the value of an unknown quantity whose probability is greatest, since that probability is always only infinitely small, than to determine the value which, if one adheres to it, makes for the least disadvantageous gamble. Or, if $fa$ denotes the probability of the value $a$ for the unknown $x$, it matters less that $fa$ attain a maximum than that $\int fx \cdot F(x-a)\,dx$, extended over all possible values of $x$, attain a minimum. Here one chooses for $F$ a function that is always positive and in a suitable way always grows larger for larger arguments. The choice of the square is purely arbitrary, and this arbitrariness lies in the nature of the matter. Were it not for the well-known, extraordinarily great advantages of choosing the square, one could choose any other function meeting those conditions, as one does in quite exceptional cases. I do not know, however, whether I have expressed myself clearly enough about what I really mean by this distinction. If not, I ask you kindly to excuse me.

. . . . . . . . . . . . . . . .

\Needspace{6\baselineskip}
\begin{center}\textsc{Gauss} to \textsc{Schumacher}. Göttingen, 25 November 1844.\end{center}

. . . . . . . . . . . .

In teaching the method of least squares orally, I usually take exactly the reverse course from the sensible one in a printed paper. I teach \textls{first} how to apply it, mixing in more or fewer of the finer devices as circumstances warrant. Only then, as far as time permits, do I take up the different justifications,
